\documentclass[a4paper]{article}
\usepackage{graphicx}
\usepackage{pdfpages}
\usepackage{hyperref}
\usepackage{etoc}
\usepackage[ngerman]{babel}
\usepackage[T1]{fontenc}

% cSpell:words Filesharing techstack Dezhong Zhuang Svelte SMUI
% Benachrichtigungsdrawer Kontaktdrawer peerjs Fileplay

\begin{document}


\begin{titlepage}
  \begin{center}
    \textbf{Ernst-Mach-Gymnasium Haar}

    \vspace{0.5cm}

    Ausbildungsabschnitt 11/2

    \vspace{1cm}

    \includegraphics[width=0.2\textwidth]{fileplay.png}

    \vspace*{1cm}

    \textbf{Fileplay}

    \vspace{0.5cm}

    Dateiübertragungswebseite mit Kontaktfunktion

    \vfill
           
    \textbf{
      \textit{Quentin Frey, Leonhard Masche, Dezhong Zhuang}
    }

  \end{center}
\end{titlepage}

\setcounter{tocdepth}{2}
\tableofcontents
\newpage

\section{Einleitung}
Unser Projekt ist eine Filesharing App (Webseite), mit der Dateien Peer-to-Peer,
also von Nutzer zu Nutzer, ohne zwischengeschalteten Server, übertragen werden
können. Die App basiert auf modernen Web-Technologien, wie z.B.
\href{https://web.dev/push-notifications-web-push-protocol/}{Web Push},
\href{https://web.dev/learn/pwa/service-workers/}{Service Workers} und
\href{https://web.dev/progressive-web-apps/}{Progressive Web Apps}. Durch die
Design-Philosophie des
\href{https://developer.mozilla.org/en-US/docs/Glossary/Progressive_Enhancement}{Progressive
      Enhancement} werden aber auch ältere Browser unterstützt. Eine Besonderheit der
App ist es, dass sie dauerhafte Verbindungen durch Kontakte möglich macht.
Kontakte können zwischen mehreren Nutzern geknüpft werden, wobei jeder Nutzer
mehrere Geräte haben kann. Zum Schutz der Privatsphäre und Benutzersicherheit
werden keine identifizierbaren Informationen auf dem Server (in der Datenbank)
gespeichert und die Dateien im Transit Ende-zu-Ende verschlüsselt. Die
Motivation für dieses Projekt war das Problem, dass es keine einfache
Möglichkeit gibt, zwischen iOS / Android / Windows / Linux Geräten Dateien zu
teilen. Bestehende Webseiten nutzen hierfür einen Link, den der Empfänger öffnen
muss, um die Datei zu empfangen. Hierbei stellt sich dann wieder die Frage, wie
man diesen Link teilt. Um dieses Problem zu behandeln, implementiert die App ein
Kontakt-System, in dem sich Nutzer einmalig über einen 6-stelligen Code
verbinden und von nun an einfach miteinander Dateien teilen können. Dieses
System ist vor allem in einer Umgebung nützlich, in der man oft Dokumente mit
denselben Personen teilt, wie es z.B. in der Schule der Fall ist.

\section{Dokumentation}

\subsection{Implementierung}
Die App wurde in Form einer Webseite realisiert und ist somit auf allen Geräten
zugänglich, die einen Browser haben. Für die Entwicklung haben wir uns für das
Framework \href{https://kit.svelte.dev/}{SvelteKit} im NodeJS Ökosystem
entschieden. Dadurch muss die App nicht manuell in den Browser-kompatiblen
Sprachen geschrieben werden, sondern kann in einer modernen Sprache (Svelte in
Verbindung mit Typescript) geschrieben werden, wodurch die Entwicklung durch
vorgefertigte Implementierungsmuster und Tools stark vereinfacht wird. Die App
ist nach dem Muster einer Single-Page-Application aufgebaut, d.h. die Seite wird
nur einmal geladen, und alle weiteren Interaktionen finden über JavaScript
statt. Dadurch wird die App sehr schnell, da nur die Daten, die benötigt werden,
vom Server geladen werden müssen. Zusätzlich wird die App dadurch sehr
benutzerfreundlich, da der Nutzer nicht durch das Laden von verschiedenen Seiten
unterbrochen wird. Die Dateiübertragung selbst findet über die Bibliothek PeerJS
statt, welche wiederum WbRTC benutzt. Durch die optionale Installation als
Progressive Web App (PWA) kann die App auch auf dem Homescreen installiert
werden, und verhält sich dann wie eine native App.\\
Die Entwicklung der App findet in einer öffentlichen GitHub Repository statt.
Dieses Repository ist unter \url{https://github.com/leonhma/fileplay} zu finden,
und wird automatisch durch Integration mit Cloudflare auf den Adressen
\url{https://app.fileplay.me} (für den {main}-branch) und
\url{https://dev.fileplay.pages.dev} (für den {dev}-branch) zur Verfügung
gestellt.\\
Der Programmcode für einen Microservice der die Kommunikation von Server zu
Client verwaltet ist unter \\\url{https://github.com/leonhma/fileplay-worker} zu
finden. Dieser Mikroservice nutzt Cloudflare APIs die nicht in SvelteKit
verfügbar sind, und wird deshalb separat gehosted.

\subsection{Architektur}
Die App ist nach dem MVC-Modell in drei Teile aufgeteilt: Frontend, Backend und
Datenbank. Das Frontend (View) ist die Webseite, die der Nutzer sieht und mit
der er interagiert. Das Backend (Controller) ist der Server, der die Webseite
ausliefert und die Anwendungslogik verwaltet, und die Datenbank ist das Modell,
in dem die Daten gespeichert werden. Sowohl das Frontend, als auch große Teile
des Backends werden vom Framework SvelteKit verwaltet. Dieses Framework erlaubt
es, die Webseite mit einer angepassten Syntax zu schreiben und
Entwickler-freundlich in Komponenten zu zerteilen. Zusätzlich bietet es eine
einfache Möglichkeit, Backend und Frontend Seite an Seite zu entwickeln.
Zusätzlich zu dem Backend-Code, der in SvelteKit geschrieben ist, gibt es noch
einen Microservice, der als \href{https://workers.cloudflare.com/}{Cloudflare
      Worker} in Typescript geschrieben ist. Dieser Verwaltet die WebSockets, die für
die Kommunikation von Server zu Client verwendet werden. Die Datenbank ist eine
SQLite Datenbank, welche mit dem Backend verbunden ist und Daten wie Geräte,
Nutzer oder Kontakte speichert.\\
Die einzelnen Teile der App (Backend, Datenbank und Microservice) werden alle
auf Cloudflare gehosted. Cloudflare bietet verschieden Web-Hosting-Services an,
darunter auch Pages, welches über eine GitHub Integration einfach mit dem
Build-Output von SvelteKit verbunden werden kann. Zusätzlich wird eine SQLite
Datenbank namens D1 verwendet, welche im Backend einfach durch
Umgebungsvariablen eingebunden werden kann. Der Microservice ist ein Cloudflare
Worker, welcher mithilfe der neuen WebSocket Hibernation API ressourcenschonend
die Verbindungen von Server zu Client verwaltet.\\
Die clientseitige Logik der App findet vor allem in Form von JavaScript
clientseitig statt. Wenn Daten vom Server abgefragt werden müssen, wie zum
Beispiel die Liste der Kontakte, wird dies über eine REST API gemacht. Dagegen
findet die Kommunikation von Server zu Client über WebSockets statt. Alternativ
wird hier auch eine Implementierung mit Web Push verwendet, jedoch wurde diese
aufgrund von Kompatibilitätsproblemen mit Edge und Firefox kurzfristig
deaktiviert. Die Dateiübertragung zwischen zwei Clients findet über WebRTC
statt. So kann die Datei (bis auf die Ausnahme einer Firewall) direkt zwischen
zwei Clients übertragen werden. Zwischen dem Front- und Backend steht noch ein
Service Worker, der im Browser der Nutzers läuft. Dieser ist für das Empfangen
von Nachrichten des Servers, und das cachen von Dateien zuständig.


\subsection{Authentifizierung}

Beim ersten Besuch der Webseite wird der Nutzer aufgefordert, bei der
Einrichtung einen Name für das Gerät und den verbundenen Nutzer einzugeben.
Diese Informationen werden dann an den Server gesendet, welcher dem Gerät eine
einzigartige Device ID zuweist. Diese ID wird dann in einem Cookie gespeichert,
und bei weiteren Besuchen der Webseite an den Server gesendet. Der Server kann
dann durch diese ID das Gerät identifizieren, und die Daten des Geräts aus der
Datenbank laden.\\
Um zu verhindern, dass ein Nutzer sich als ein anderes Gerät ausgibt, wird
zusätzlich zur Device ID ein Token generiert, welches auch einem Cookie
gespeichert wird. Dieses Token wird bei jeder Anfrage an den Server gesendet,
und der Server kann dann überprüfen, ob das Token, das mit der Signatur des
Servers unterzeichnet wurde, zu der Device ID passt.

\subsection{Geräte, Nutzer, Kontakte}

Um eine einfache Organisation und Nutzung der App zu ermöglichen, wird zwischen
Geräten, Nutzern und Kontakten unterschieden. Ein Gerät ist ein Endgerät, auf
dem die App installiert ist. Ein Nutzer ist eine Person, die die App nutzt. Ein
Nutzer kann mehrere Geräte besitzen, ein Gerät dagegen kann nur einem Nutzer
zugeordnet werden. Ein Kontakt repräsentiert eine Verbindung zwischen zwei
Nutzern, welche wiederum unbegrenzt viele Kontakte haben können.\\
Diese Verknüpfungen von Nutzern und Geräten (one-to-many) und Nutzern
untereinander, als Kontakte (many-to-many), können durch das Generieren und
Einlösen von sechsstelligen Verbindungscodes realisiert werden. Wenn ein Nutzer
einen Code generiert, wird dieser Code in der Datenbank gespeichert und kann
diesen Code an einen anderen Nutzer weitergeben. Wenn der andere Nutzer diesen
Code eingibt, wird eine Verbindung zwischen den beiden Nutzern hergestellt und
der Code wird aus der Datenbank gelöscht.


\subsection{Filesharing}
Wenn ein Nutzer eine Datei teilen möchte, wählt er einen verfügbaren Kontakt
aus, dessen ID dann in einer Anfrage an den Server gesendet wird. Dieser schickt
diese Anfrage anschließend per WebSocket (oder Web Push) an die einzelnen Geräte
des Kontakts. Wenn ein Gerät die Anfrage akzeptiert, schickt dieses Gerät eine
Antwort mit der PeerJS-Adresse des Geräts, an das die Datei übertragen werden
soll, an den Server. Diese Antwort enthält auch den Public Key des Geräts, der
für die Verschlüsselung der Datei verwendet wird. Der Server leitet diese
Antwort dann an den ursprünglichen Nutzer weiter, der die Datei dann an die
PeerJS-Adresse senden kann.

\subsection{User-Story}

\subsubsection{Einrichtung}
Wenn ein Nutzer die App zum ersten Mal öffnet, wird er aufgefordert, einen Namen
für das Gerät anzugeben, und den Typ des Geräts auszuwählen. Dies dient
lediglich der erleichterten späteren Identifikation der Geräte durch den Nutzer
selbst.\\
Zusätzlich zu den Daten des Geräts muss der Nutzer entweder einen neuen Nutzer
erstellen, oder das Gerät zu einem Bereits existierenden hinzufügen. Wenn der
Nutzer einen neuen erstellt, Kann er diesem einen Namen geben. Dieser Name wird
auch für Kontakte sichtbar sein, wenn mit ihnen eine Datei geteilt wird.
Dementsprechend wird der Name mithilfe der kostenlosen API \url{
      https://www.purgomalum.com/} auf unangemessene Inhalte überprüft.\\
Ein weiteres identifizierendes Merkmal des Nutzers ist ein Avatar. Um auch hier
ungewollte Überraschungen zu vermeiden, wird lediglich ein Seed generiert, der
dann durch die \href{ https://www.dicebear.com/}{Dicebear} API in einen Avatar
umgewandelt wird. Wenn der Nutzer das Gerät zu einem bereits existierenden
Nutzer hinzufügt, muss er den Verbindungscode eingeben, den der Nutzer generiert
hat. Dieser Code wird dann an den Server gesendet, der dann überprüft, ob der
Code existiert. Wenn der Code existiert, wird das Gerät dem Nutzer hinzugefügt,
und der Code wird aus der Datenbank gelöscht. Wenn der Code nicht existiert,
wird eine Fehlermeldung angezeigt.\\
Nach der erfolgreichen Einrichtung wird sich die Statusanzeige (in der oberen
rechten Ecke) schnell auf grün ändern. Das bedeutet, dass das Gerät mit dem
Server verbunden ist, und die App vollständig einsatzbereit ist.

\subsubsection{Kontakte verwalten}

Zum hinzufügen und entfernen von Kontakte muss sich der Nutzer zuerst die
Einrichtung durchlaufen. Dann kann er in der Top-Bar (der orangen Leiste oben im
Fenster) auf das Icon mit dem Kontaktbuch klicken. Daraufhin öffnet sich ein
Drawer (eine Leiste, die von der Seite aus aufgezogen wird), in dem alle
Kontakte aufgelistet sind. Wenn der Nutzer einen Kontakt hinzufügen möchte, kann
er auf 'add contact' klicken. Daraufhin öffnet sich ein Dialog, in dem der
Nutzer entweder einen bestehenden Verbindungscode eingeben, oder einen neuen
generieren kann. Wenn der Nutzer einen neuen Code generiert, wird dieser Code in
der Datenbank gespeichert, und der Nutzer kann diesen Code an einen anderen
Nutzer weitergeben.\\
Dieser Andere Nutzer kann dann in demselben Dialog auf seinem Gerät den Code eingeben,
und verbindet sich somit als Kontakt der Person, die den Code generiert hat.
Diese Verbindung ist beidseitiger Natur, d.h. beide Nutzer sind Kontakte des
jeweils anderen.\\
Möchte der Nutzer einen Kontakt entfernen, kann er in der Liste der Kontakte auf
das Trash-Icon neben dem Kontakt klicken. Daraufhin wird eine Anfrage an den
Server gesendet, der dann den Kontakt aus der Datenbank löscht. Diese Anfrage
bedarf natürlich - so wie die meisten - Authentifizierung damit nicht jeder
Nutzer die Kontakte anderer Nutzer löschen kann.



\subsubsection{Dateien teilen}

Um eine Datei zu teilen, muss der Nutzer zuerst die Einrichtung durchlaufen.
Dann kann er mitten auf der Seite das Icon mit dem Upload-Symbol klicken.
Daraufhin wird vom Browser ein Datei-Dialog geöffnet, in dem der Nutzer eine
Datei auswählen kann. Wenn der Nutzer eine Datei ausgewählt hat, wird der Name
der Datei neben dem Icon angezeigt. Nun kann der Nutzer durch einen Klick auf
das 'Senden' Icon die Datei an einen Kontakt senden. Es öffnet sich ein Dialog,
in dem der Nutzer einen Kontakt auswählen kann. Hier werden nur Kontakte
angezeigt, die in der letzten Minute ein Lebenszeichen an den Server gesendet
haben. Wenn der Nutzer einen Kontakt auswählt, wird eine Anfrage an den Server
gesendet, der dann die Anfrage an die einzelnen Geräte des Kontakts weiterleitet
und auf eine Antwort wartet. Während diesem Prozess steht der Status beim
Sender auf 'requesting'. Wenn ein Gerät die Anfrage akzeptiert, wird ein
weiterer Request an den Server gesendet, der auch die Details benötigt, um die
Datei zu senden. Wenn der Server diese Details erhalten hat, wird eine
Benachrichtigung mit den Details an das Gerät gesendet, welches die
ursprüngliche Anfrage gestartet hat. Dieses Gerät kann dann die Datei an das
Gerät des Kontakts senden. Während der Übertragung der Datei wird der Status
beim Sender auf 'sending' gesetzt. Passiert ein Fehler, wird der Status auf
'failed' gesetzt.

\subsubsection{Nutzer und Geräte verwalten}

Wenn der Nutzer die Details (Name und Avatar) seines Nutzers ändern möchte, kann
er in der Top-Bar auf das Icon mit dem Zahnrad klicken, und zum Tab 'Account'
navigieren. Daraufhin öffnet sich ein Dialog, in dem der Nutzer seinen Namen und
Avatar ändern kann. Wenn der Nutzer seinen Namen oder Avatar ändert, wird dieser
auch für alle Kontakte geändert. Wenn der Nutzer die Namen oder Typen seiner
Geräte ändern möchte, so kann er das im Tab 'Devices' tun. Hier werden alle
Geräte des Nutzers aufgelistet, und durch einen Klick auf das Dreipunkt-Menü wird
ein weiterer Dialog geöffnet, in dem die Details des Geräts angepasst,
sowie die Verbindung des Gerätes zum Account gelöst werden können.

\subsubsection{App-Updates}

Während der Entwicklung, aber auch während der Veröffentlichung der App, ist es nützlich,
den Nutzer über Updates zu informieren. Das ist besonders wichtig, da ein
Service Worker benutzt wird, der die Webseite offline verfügbar mach, und somit
nicht immer die aktuellste Version der Webseite lädt. Um den Nutzer über Updates
zu informieren, wird periodisch (30 Minuten) überprüft, ob eine neue Version der
App verfügbar ist. Ist das der Fall, wird eine Benachrichtigung angezeigt, die
den Nutzer informiert, dass ein Update verfügbar ist. Wenn der Nutzer auf diese
Benachrichtigung klickt, wird die Seite neu geladen, und die neue Version wird
installiert.\\
Durch die Integration mit GitHub und Cloudflare wird die App automatisch bei
jedem Commit auf GitHub neu gebaut und auf Cloudflare Pages veröffentlicht. Die
Entwicklung findet auf dem 'dev'-branch statt, und finalen Versionen werden über
einen Pull-Request auf den 'main'-branch übertragen. Wenn ein Pull-Request auf
den 'main'-branch gemerged wird, wird die App automatisch auf der Haupt-Domain
'app.fileplay.me' veröffentlicht.


\section{Aufteilung der Aufgaben}

\addto{\captionsngerman}{\renewcommand*{\contentsname}{Mitglieder}}
\setcounter{tocdepth}{2}
\localtableofcontents
\addto{\captionsngerman}{\renewcommand*{\contentsname}{Inhaltsverzeichnis}}

\subsection{Leonhard Masche (\href{https://github.com/leonhma}{@leonhma})}
\begin{itemize}
      \item Infrastruktur und DevOps: Ich habe die Struktur des Projektes
            eingerichtet, und die Infrastruktur auf Cloudflare konfiguriert.
            Dazu gehört die Konfiguration des Domainnamen und DNS, sowie das
            Einrichten der Datenbank, deren Integration mit dem Backend, sowie
            die CI/CD Integration mit GitHub. Implementierungen sind:
            '/src/lib/db.ts', '/src/lib/server/db.ts', '/wrangler.toml',
            '/migrations/*'
      \item API: Ich habe die REST API entwickelt, die z.B. das Auflisten von
            Kontakten und Geräten, Anfragen von Datei-Sharing und das Verbinden
            von Nutzern untereinander und mit Geräten ermöglicht. Im Projekt
            sind all diese Dateien unter '/src/routes/api/**' zu finden.
      \item Authentifizierung: Ich habe die Authentifizierung implementiert, die
            dafür sorgt, dass nur authentifizierte Geräte mit dem Server
            kommunizieren können. Dazu gehört die Generierung von Tokens und
            Device-IDs, sowie die Überprüfung dieser. Die Implementation ist
            unter '/src/lib/server/crypto.ts' zu finden.
      \item Server Side Notifications: Schon beim Thema Backend angekommen, habe
            ich auch das System für die Benachrichtigungen konzeptioniert und
            implementiert. Dieses System besteht aus einer Funktion auf dem
            Server, die eine Nachricht adaptiv über einen der beiden Kanäle
            (WebSockets oder Web Push) an den Client überliefert, sowie der
            clientseitige Logik, um die Verbindung für die Nachrichten zu
            initialisieren und aktiv zu halten. Die Implementation befindet sich
            unter '/src/lib/server/notifications.ts' udn '/src/sw.ts'
      \item Web Push: Als einer der beiden Kanäle für die Benachrichtigungen
            habe ich Web Push implementiert. Dazu gehört auch die
            Authentifizierung mit dem VAPID-Keypair des Servers. (siehe
            '/scripts/generate-keys.js'). Dieser Kanal ist vor allem für die
            Benachrichtigungen auf mobilen Geräten interessant, da er
            stromsparend Benachrichtigungen über einen Zentralisierten
            Push-Service liefert, auch wenn das Gerät sich nicht auf der
            Webseite befindet. Die Implementation findet sich zum Teil in
            '/src/lib/server/notifications.ts' und zum Teil im Service Worker
            ('/src/sw.ts'). Aufgrund von unergründbaren Problemen bei der
            Kompatibilität mit anderen Browsern (nicht Google Chrome), musste
            dieses Feature (zumindest für die Vorführung) aber deaktiviert
            werden.
      \item Service Worker: Zwischen Client und Server habe ich auch den Service
            Worker geschrieben, der das Caching der Webseite für eine schnellere
            Ladezeit mithilfe WorkBox verwaltet, und eine wichtige Rolle bei den
            Push Notifications spielt. Der Service Worker befindet sich unter
            '/src/sw.ts'.
      \item WebSockets: Als zweiter Kanal für die Benachrichtigungen habe ich
            die finale Version der WebSockets implementiert. Zu dieser
            Benachrichtigungsmethode gehört auch der in
            'leonhma/fileplay-worker' aufzufindende Cloudflare Worker, der ein
            Durable Object verwaltet, welches die WebSocket Verbindung dem
            Server gegenüber zugänglich macht.
      \item Als ein weiterer Teil der Notifications, habe ich mich mit der
            Übertragung von Nachrichten zwischen Service Worker und Client
            auseinandergesetzt. Dazu gehört auch, die per Web Push empfangenen
            Nachrichten an den Client (die Webseite) weiterzuleiten, um auf sie
            reagieren zu können. Die relevanten Stücke Code befinden sich in
            '/src/sw.ts' und /src/lib/messages.ts'.
      \item Notifications: Als letzter Teil der Notifications habe ich mich um
            die Anzeige von interagierbaren Benachrichtigungen gekümmert. Das
            sind zum einen die Benachrichtigungen die über den Service Worker
            und Web Push im OS-Tray angezeigt werden, aber auch die Nachrichten,
            die in der App im Notification Drawer sichtbar sind. Bei den In-App
            Benachrichtigungen habe ich mich um die klickbaren Aktionen und
            deren Ausführung bemüht, welche stark die Notification API im
            Service Worker nachahmen. Die relevanten Dateien sind '/src/sw.ts',
            '/src/lib/stores/Dialog.ts' und
            '/src/lib/components/NotificationDrawer.svelte'.
      \item Ein Teil des UIs den Ich gestaltet habe, ist der
            Einrichtungs-Dialog. Dieser ist eng Verbunden mit der
            Authentifizierung und Der REST API. Die relevante Datei ist
            '/src/lib/dialogs/SetupDialog.svelte'.
\end{itemize}

\subsection{Quentin Frey (\href{https://github.com/limosine}{@limosine})}
\begin{itemize}
      \item UI (Material Design):\\
            Viele Webseiten, Progressive Web Apps oder auch native Apps
            verwenden Material Design. Um die Eingewöhnung so einfach wie
            möglich zu machen, hielt ich mich streng an die Material Design
            Guidelines. Daher verwendete ich ausschließlich Komponenten aus SMUI
            (SvelteMaterialUI), zudem blieb der Code dadurch übersichtlich. Die
            UI der App besteht aus der Startseite, dem Benachrichtigungsdrawer,
            dem Kontaktdrawer und einer Leiste, welche sich am oberen
            Fensterrand befindet. Außerdem können einige Dialoge, wie der
            Einstellungsdialog, der Dialog für das Senden von Dateien oder der
            Dialog, welcher es ermöglicht die Eigenschaften von Geräten zu
            verändern, geöffnet werden. In der oberen App-Leiste befinden sich
            Knöpfe, mit welchen die Drawer geöffnet werden können, der aktuelle
            Status angezeigt wird und der Einstellungsdialog geöffnet werden
            kann.
      \item State Management:\\
            Da die angezeigten Daten immer auf dem aktuellsten Stand sein
            müssen, wurde ein zentraler Timer implementiert, welcher die Daten
            in einem bestimmten Intervall abruft, wenn diese benötigt werden.
            Damit die Daten nicht verloren gehen und die Seite dynamisch
            erneuert wird, werden Svelte Stores verwendet, in denen die oben
            genannten Daten gespeichert werden.
      \item Senden \& Aufteilen von Dateien (PeerJS):\\
            Die Dateien werden in 1 Megabyte große Stücke (Chunks) aufgeteilt
            und mithilfe von PeerJS, welches die Nutzung von WebRTC-Protokollen
            vereinfacht, versendet. Dies ist im Ordner \textit{src/lib/peerjs}
            zu finden.
      \item Verschlüsselung \& Entschlüsselung von Dateien (OpenPGP):\\
            Da die OpenPGP-Verschlüsselung schon im Alltag verwendet wird, ist
            diese vielfach überprüft worden und daher sehr sicher.
            Beispielsweise beim Versenden von E-Mails wird oftmals die
            Verschlüsselung durch GPG / OpenPGP durchgeführt. Bei diesem
            Verfahren wird ein öffentlicher und ein privater Schlüssel erstellt.
            Um eine Nachricht zu verschlüsseln, wird der öffentliche Schüssel
            des Empfängers benötigt, sodass der Empfänger die Datei mit seinem
            persönlichen privaten Schlüssel entschlüsseln kann. Diese
            Nachrichten werden durch eine Armor verstärkt, indem diese ins
            ASCII-Encoding konvertiert werden, anschließend werden diese an den
            Empfänger gesendet, welcher die Stücke dann wieder zusammensetzt und
            entschlüsselt.
      \item WebSockets (ohne Durable Objects):\\
            Durable Objects sind Objekte, welche im Cloudflare-Netzwerk in jedem
            Rechenzentrum existieren und gegenseitig miteinander vernetzt sind.
            Durch diese Objekte ist es möglich, mehrere WebSocket-Verbindungen
            zu koppeln, sodass Daten ohne Verzögerung übertragen werden können.
            Da Durable Objects Premium-Funktionen von Cloudflare sind,
            implementierte ich eine WebSocket-Verbindung, welche nicht auf
            Durable Objects oder ähnliches angewiesen ist. Diese wurde aber von
            dem Fileplay-Worker ersetzt, sodass die Webseite, nachdem das
            Abonnement gekündigt wurde, nicht mehr funktioniert und
            voraussichtlich durch die ältere Methode ersetzt werden muss. Diese
            ältere Methode verwendete Einträge in der Datenbank, um die
            PeerJS-ID sowie den öffentlichen Schlüssel zu übertragen und
            funktioniert daher auch, wenn die WebSockets keine gegenseitige
            Kenntnis voneinander haben.
      \item Gast-Modus (ohne Erstellung eines Accounts):\\
            Da viele Personen lediglich ein einziges Mal eine Datei empfangen
            möchten, ist es für diese umständlich einen Account erstellen zu
            müssen. Daher fügte ich den Gast-Modus hinzu, welcher es ermöglicht
            einen Link zu generieren, welchen jeder aufrufen kann, ohne einen
            Account zu erstellen. Auch bei diesem Modus wird die Datei
            verschlüsselt, statt einem PublicKey wird hier jedoch einen
            Zeichenfolge verwendet, welche im Link enthalten ist.
\end{itemize}

\subsection{Dezhong Zhuang (\href{https://github.com/didi1150}{@didi1150})}
\begin{itemize}
      \item BETA UI:\\
            Für das Design einer Website/App muss zuerst immer ein Design
            entworfen werden, welches als Vorlage für die zukünftige Entwicklung
            dient. Dazu wurde explizit pures HTML, CSS und Javascript verwendet,
            um den Designprozess so schnell und effizient wie möglich zu
            gestalten. Zudem führt das Fehlen von jeglichen Abhängigkeiten dazu,
            dass später Features schneller implementiert werden können.
      \item Filesharing Progress UI:\\
            Der Zweck des Aufteilens der Dateien in Chunks war, dass somit der
            Empfänger Rückmeldung bekommt, wie weit der Dateiempfang schon
            vorangeschritten ist. Hierbei werden die angekommenen Chunks
            innerhalb eines Svelte Stores gespeichert (mit jeweils
            unterschiedlichen IDs zur Zuordnung) dessen Daten dann mit der Seite
            geteilt werden.
      \item State Management:\\
            Für das Suchen oder auch Speichern von ausgewählten Elementen werden
            Svelte Stores verwendet, um automatische interne Benachrichtigungen
            beim Ändern eines Wertes zu verschicken. Für den Fall, dass keine
      \item Senden \& Aufteilen von Dateien (PeerJS):\\
            Der Sender und Empfänger jeweils verschicken jeweils Metadaten
            mithilfe Peerjs aneinander, wodurch ein Transferprozess
            initialisiert wird. Während dieses Prozess wird die angefragte Datei
            in 1 Megabyte große Stücke (Chunks) aufgeteilt und an den Empfänger
            gesendet, Während dieser nach jedem Chunk eine Bestätigung sendet.
      \item Recherche und Vorstrukturierung:\\
            Für die Verwendung von neuen Resourcen muss oft vorrecherhiert
            werden, welche APIs oder SDKs nutzbar sind oder den Anforderungen
            des Projekts entsprechen. Im Anschluss daran werden diese
            vorübergehend in ihrer gewünschten Funktion zum Testen
            implementiert, damit diese später optimiert werden können.
      \item Generierung eines QR-Codes:\\
            Während des Testens mit mehr oder weniger freiwilligen Probanden
            wurde festgestellt, dass es sehr umständlich ist, einen Link an
            andere Personen ohne Verwendung von externen Plattformen zu
            verschicken. Aus diesem Grund wurde die Anzeige eines generierten
            QR-Codes implementiert, welcher einfach und leicht aufzurufen ist.
\end{itemize}

\end{document}